\subsection{Переход от OD-пар к источникам}

Теоретически блоком может быть отдельная OD-пара $w=(i,j)$.
Однако в дорожном графе эффективнее выбирать источники.
Если выбран источник $i$, то один запуск алгоритма кратчайших путей из $i$ дает расстояния
до всех стоков $j$.
Поэтому батч источников
\[
    \mathcal A_k\subset O
\]
порождает батч OD-пар
\[
    \cI_k=\{(i,j)\in OD\mid i\in\mathcal A_k\}.
\]
Такой выбор лучше соответствует структуре вычислений:
стоимость одной итерации примерно пропорциональна числу выбранных источников, а не числу
отдельных OD-пар.

\subsection{Обновление разложения потока}

Пусть
\[
    f_k=\sum_{i\in O}f_k^{(i)}.
\]
Для выбранного источника $i$ строится новый all-or-nothing поток $\hat f_k^{(i)}$:
весь спрос из $i$ назначается на кратчайшие пути при текущих временах $\tau(f_k)$.
Затем направление для этого источника равно
\[
    d_k^{(i)}=\hat f_k^{(i)}-f_k^{(i)}.
\]
Если выбран батч $\mathcal A_k$, суммарное направление записывается как
\[
    d_k=\sum_{i\in\mathcal A_k}d_k^{(i)}.
\]
Обновление выполняется только для выбранных источников:
\[
    f_{k+1}^{(i)}
    =
    \begin{cases}
        f_k^{(i)}+\gamma_k d_k^{(i)}, & i\in\mathcal A_k,\\
        f_k^{(i)}, & i\notin\mathcal A_k.
    \end{cases}
\]
После этого
\[
    f_{k+1}=\sum_{i\in O}f_{k+1}^{(i)}.
\]
Эта формула показывает, почему SOFW требует хранения $f^{(i)}$.
Без такого разложения нельзя локально заменить вклад выбранного источника, не пересчитывая
полный поток.

\subsection{Равномерная и взвешенная выборка}

При равномерной выборке каждый источник имеет одинаковую вероятность попасть в батч.
Это просто и дает чистую теоретическую оценку.
Но спрос по источникам может быть сильно неравномерным.
Если источник генерирует большой поток, его обновление может сильнее уменьшить функционал.
Поэтому рассматривается взвешенная версия.

Пусть
\[
    D_i=\sum_j d_{ij}
\]
-- полный спрос из источника $i$.
Вероятность выбора источника задается как
\[
    p_i=\frac{D_i}{\sum_{\ell\in O}D_\ell}.
\]
Чтобы направление не стало смещенным в сторону крупных источников, можно использовать
importance weighting.
В простейшей записи вклад выбранного источника корректируется множителем $1/p_i$:
\[
    \tilde d_k^{(i)}=\frac{1}{p_i}d_k^{(i)}.
\]
Тогда
\[
    \E[\tilde d_k^{(i)}]=\sum_{i\in O}d_k^{(i)}
\]
с точностью до нормировки батча.
На практике эта коррекция стабилизирует сравнение равномерного и взвешенного SOFW.

\subsection{Влияние батча на теоретический порядок}

Интуитивно может показаться, что батч размера $m$ должен ускорять сходимость в $m$ раз.
Однако в нормированной записи направление усредняется по выбранным блокам.
Ожидаемый прогресс на одну итерацию остается пропорционален полному FW-зазору, деленному
на число блоков:
\[
    \E[\Delta_k\mid x_k]\approx -\frac{\gamma_k}{|OD|}g(x_k).
\]
Батч уменьшает дисперсию, но не меняет этот множитель в базовой оценке.
Практический выигрыш появляется через время выполнения:
если батч можно обработать значительно быстрее полного оракула, то за тот же бюджет времени
алгоритм делает больше обновлений.

\subsection{Память и время}

Если в графе $|O|$ источников и $|\cE|$ ребер, хранение полного потока требует
\[
    \cO(|\cE|)
\]
памяти.
Хранение разложения по источникам требует
\[
    \cO(|O||\cE|)
\]
в худшем случае.
На практике потоки могут быть разреженными, но порядок показывает потенциальную цену SOFW.
Время одной полной итерации FW можно оценить как
\[
    |O|\cdot T_{\mathrm{SP}},
\]
где $T_{\mathrm{SP}}$ -- стоимость одного запуска кратчайших путей.
Если SOFW выбирает долю $\alpha$ источников, то итерация стоит примерно
\[
    \alpha |O|\cdot T_{\mathrm{SP}},
\]
но требует поддерживать разложение $f=\sum_i f^{(i)}$.

\subsection{Итоги раздела}

Стохастический блочно-координатный подход показывает, что для крупных транспортных сетей главная стоимость FW -- это
не формула шага и не вычисление градиента, а полный линейный оракул.
SOFW уменьшает эту стоимость, выбирая только часть источников.
Теоретическая оценка остается такой же по порядку, как у блочного FW, а практические графики
показывают, что на больших сетях экономия времени важнее шума стохастического направления.
